\documentclass{article}
\usepackage{amsmath}
\usepackage{xcolor}
\begin{document}

VARIACION DE PARAMETROS

- Este metodo es generico para cualquier tipo de E.D. con CC. no Homogenea

Permite determinar la mayor parte de soluciones de un E.D. 

Depende de la solucion homogenea

Requiere el Wronskiano como metodo de sistema de ecuaciones

Se apoya en el uso de Derivadas e Integrales.

Dada una E.D. de la forma:




\begin{gather*}
a_{1}(x)y''+a_{2}(x)y'+a_{3}(x)y = g(x) \\
\vspace{0.5em} \\
1)\text{ Determinar la solucion Homogenea } \rightarrow  y_{h} \\
2)\text{ Poner la ecuacion Diferencial en su forma estandar} \\
y''+\frac{a_{2}}{a_{1}(x)}(x)y'+\frac{a_{3}}{a_{1}(x)}(x)y =\frac{g(x)}{a_{1}(x)} \\
\vspace{0.5em} \\
3)\text{ Se plantea una solucion generica, en base a las solucion y}_{h} \\
\text{por decir:} \\
y_{h} = c_{1}y_{1}+c_{2}y_{2} \\
\vspace{0.5em} \\
\text{se asume una solucion generica en base a y}_{h} \\
y1 = x^{2}  y2 = x^{3}   \rightarrow  v1=1, v2 =\frac{1}{x} \\
y_{p} = v_{1}y_{1}+v_{2}y_{2} =   \rightarrow \text{ donde v}_{1}, y v_{2}  \rightarrow \text{funciones de x desconocidas} \\
\text{el objetivo es determinar v1, y v2} \\
\vspace{0.5em} \\
4)\text{ Para determinar v1 y v2, se arma un sistema de ecuaciones,} \\
\text{similar al Wronskiano} \\
v_{1}'y_{1}+v_{2}'y_{2} = 0 \\
v_{1}'y_{1}'+v_{2}'y_{2}' =\frac{g(x)}{a_{1}(x)} \\
\vspace{0.5em} \\
\text{si el sistema fuese de 4to orden} \\
v_{1}'y_{1}+v_{2}'y_{2}+v_{3}'y_{3}+v_{4}'y_{4} = 0 \\
v_{1}'y_{1}'+v_{2}'y_{2}'+v_{3}'y_{3}'+v_{4}'y_{4}' = 0 \\
v_{1}'y_{1}''+v_{2}'y_{2}''+v_{3}'y_{3}''+v_{4}'y_{4}'' = 0 \\
v_{1}'y_{1}'''+v_{2}'y_{2}'''+v_{3}'y_{3}'''+v_{4}'y_{4}''' = \frac{g(x)}{a_{1}(x)} \\
\vspace{0.5em} \\
\matrix{y_{1}}{y_{1}'}{y_{1}''}{y_{1}'''}{y_{2}}{y_{3}}{y_{4}}{y_{2}'}{y_{3}'}{y_{4}'}{y_{2}''}{y_{3}''}{y_{4}''}{y_{2}'''}{y_{3}'''}{y_{4}'''}\matrix{v_{1}'}{v_{2}'}{v_{3}'}{v_{4}'}=\matrix{0}{0}{0}{\frac{g(x)}{a_{1}(x)}} \\
6)\text{ Determinar el valor v}_{1}',v_{2}'....v_{n}' \\
7)\text{ Determinamos las integrales respectiva} \\
v_{1}' = f(x)   \rightarrow  v_{1} = \int f(x)dx \\
8)\text{ Reemplzar en la funcion supuesta} \\
9)\text{ Reducir algebraicamente en caso de ser posible} \\
10)\text{ Absorver los terminos similares en caso de ser posible} \\
11)\text{ Determinar la solucion general como tal} \\
y = y_{h}+y_{p}
\end{gather*}


Ejemplos




\begin{gather*}
y′′′−2y′′−y′+2y=e^{3x} \\
e^{rx}(r^{3}-2r^{2}-r+2) = 0 \\
\matrix{1}{-2}{}{1}{1}{-1}{-1}{2}{}{-1}{-2}{1}{-2}{0}{} \\
(r^{2}-r-2)(r-1)=0 \\
(r-2)(r+1)(r-1)=0 \\
yh= c_{1}e^{2x}+c_{2}e^{x}+c_{3}e^{-x} \\
y_{p} = Ae^{3x} \\
y_{p} = \frac{1}{8}e^{3x} \\
\vspace{0.5em} \\
y_{h}= v_{1}e^{2x}+v_{2}e^{x}+v_{3}e^{-x} \\
\text{Armamos el sistema de Ecuaciones} \\
 v_{1}'\text{\colorbox[RGB]{248,231,28}{$e$}}\text{\colorbox[RGB]{248,231,28}{$^{2x}$}}+v_{2}'\text{\colorbox[RGB]{248,231,28}{$e$}}\text{\colorbox[RGB]{248,231,28}{$^{x}$}}+v_{3}'\text{\colorbox[RGB]{248,231,28}{$e$}}\text{\colorbox[RGB]{248,231,28}{$^{-x}$}}=0    \rightarrow (A) \\
 2v_{1}'e^{2x}+v_{2}'e^{x}-v_{3}'e^{-x}=0  \rightarrow (B) \\
4v_{1}'e^{2x}+v_{2}'e^{x}+v_{3}'e^{-x}=e^{3x} \rightarrow (C)
\end{gather*}



\begin{gather*}
 \matrix{e^{2x}}{2e^{2x}}{4e^{2x}}{y_{1}'''}{e^{x}}{e^{-x}}{y_{4}}{e^{x}}{-e^{-x}}{y_{4}'}{e^{x}}{e^{-x}}{y_{4}''}{y_{2}'''}{y_{3}'''}{y_{4}'''}\matrix{v_{1}'}{v_{2}'}{v_{3}'}{v_{4}'}=\matrix{0}{0}{e^{3x}}{\frac{g(x)}{a_{1}(x)}}
\end{gather*}



\begin{gather*}
\text{Sumamos }(A)\text{ con }(B) \\
v_{1}'\text{\colorbox[RGB]{248,231,28}{$e$}}\text{\colorbox[RGB]{248,231,28}{$^{2x}$}}+v_{2}'\text{\colorbox[RGB]{248,231,28}{$e$}}\text{\colorbox[RGB]{248,231,28}{$^{x}$}}+v_{3}'\text{\colorbox[RGB]{248,231,28}{$e$}}\text{\colorbox[RGB]{248,231,28}{$^{-x}$}}=0    \rightarrow (A) \\
 2v_{1}'e^{2x}+v_{2}'e^{x}-v_{3}'e^{-x}=0  \rightarrow (B) \\
3v_{1}'e^{2x}+2v_{2}'e^{x} = 0  \rightarrow (D) \\
\text{Sumamos }(B)\text{ con }(C) \\
 2v_{1}'e^{2x}+v_{2}'e^{x}-v_{3}'e^{-x}=0  \rightarrow (B) \\
4v_{1}'e^{2x}+v_{2}'e^{x}+v_{3}'e^{-x}=e^{3x} \rightarrow (C) \\
6v_{1}'e^{2x}+2v_{2}'e^{x} = e^{3x}   \rightarrow (E) \\
\vspace{0.5em} \\
(E) -(D) \\
6v_{1}'e^{2x}+2v_{2}'e^{x} = e^{3x} \\
-3v_{1}'e^{2x}-2v_{2}'e^{x} = 0  \\
3v_{1}'e^{2x} = e^{3x} \\
v_{1}' = \frac{e^{3x}}{3e^{2x}} = \frac{1}{3}e^{x} \\
v_{1}= \int \frac{1}{3}e^{x}dx \\
v_{1}=\frac{1}{3}e^{x} \\
\text{Reemplazando en 3v}_{1}'e^{2x}+2v_{2}'e^{x} = 0  \rightarrow (D) \\
\vspace{0.5em} \\
3\frac{1}{3}e^{x}e^{2x}+2v_{2}'e^{x} = 0  \\
e^{3x}+2v_{2}'e^{x} = 0  \\
2v_{2}'e^{x}=-e^{3x} \\
v_{2}'=-\frac{1}{2}e^{2x} \\
v_{2} = -\frac{1}{2}\int e^{2x}dx \implies -\frac{1}{2}\frac{1}{2}e^{2x}=-\frac{1}{4}e^{2x} \\
\vspace{0.5em} \\
\text{Reemplazando en }(A) \\
 \frac{1}{3}e^{x}e^{2x}-\frac{1}{2}e^{2x}e^{x}+v_{3}'e^{-x}=0    \rightarrow (A) \\
\frac{1}{3}e^{3x}-\frac{1}{2}e^{3x}+v_{3}'e^{-x}=0 \\
-\frac{1}{6}e^{3x}=-v_{3}'e^{-x} \\
v_{3}'=\frac{1}{6}e^{4x} \rightarrow v_{3} =\frac{1}{24}e^{4x} \\
\vspace{0.5em} \\
y_{h}= \frac{1}{3}e^{x}e^{2x}-\frac{1}{4}e^{2x}e^{x}+\frac{1}{24}e^{4x}e^{-x} \\
y_{h}=e^{3x}(\frac{1}{3}-\frac{1}{4}+\frac{1}{24})=e^{3x}(\frac{8-6+1}{24})=\frac{3}{24}e^{3x} = \frac{1}{8}e^{3x}  \\
\vspace{0.5em} \\
\end{gather*}





\end{document}
